\documentclass[12pt]{article}

\usepackage[a4paper, total={6in, 8in}]{geometry}
\usepackage{textcomp}

\begin{document}
\setlength{\parindent}{0pt}

\def \t {\textrightarrow}
\def \v {\vspace{1em}}
\def \bi {\begin{itemize}}
\def \ei {\end{itemize}}
\def \s[#1] {\section*{#1}}
\def \ss[#1] {\subsection*{#1}}
\def \sss[#1] {\subsubsection*{#1}}

\s[Dopo guerra nel mondo]
\ss[Carta Atlantica]
I trattati vengono sanciti gia prima della fine \\
Il primo di questi avvicinamenti è la carta Atlantica nel 41 tra USA e UK \t documento che sancisce vicinanza ideale \t due paesi democratici che si impegnano a combattere ogni forma di fascismo \\
È un docuento di vicinanza ideale \\
I grandi protagonisti dei trattati sono Churcill, Roosevelt e Stalin \t la carta atlantica riguarda ancora l'occidente, non ingloba ancora la URSS che non puo sostenre i principi democratici che la carta sancisce \\
La carta atlantica non ingloba anche Stalin

\ss[Conferenza di Teheran]
Nel 1943 si incontrano Stalin, Roosevelt e Churcill \t per prima volta l URSS viene chiamata per delineare il futuro assetto del mondo \\
A Teheran si decide lo sbarco in Normandia \t ovvero l apertura di nuovo fronte di francia, e:
\bi
\item Roosevelt vuole grande organizzazione mondiale (che sara l ONU) davanti alla quale si possono risolvere controverse diplomaticamente
\item Stalin chiede invece rassicurazioni \t vuole che paesi vicini a URSS diventino suoi alleati \t vuole proteggere i confini dell'URSS
\item Churcill si fa promotore di dividere la germania in zone di influenza, e qua cerchera di definire queste zone di influenza per cui si incontrera nel 44 a Mosca con Stalin
\ei

\ss[Conferenza di Ialta]
Ialta è in Crimea, e si tiene nel 1945 \\
Ratifica tutti gli accordi precedenti \\
Qua l'URSS si impegna ad entrare in guerra contro giappone, e si prendono decisioni:
\bi
  \item germania deve essere divisa in 4 zone controllate da URSS, USA, UK e Francia
  \item viene deciso lo scioglimento dell'esercito tedesco
  \item denazificazione del paese
  \item criminali di guerra devono essere giudicati (in tribunale)
  \item deciso il pagamento dei danni di guerra da parte della germania
\ei
A Ialta viene dichiarato il principio atlantico \t verra smetntito da degli accordi segreti \t ma è il principio di determinazione dei popoli (si afferma questo) \\
Ialta è simbolo della fine della guerra \t ed è considerata il momenton in cui il mondo si divide in due zone di influenza, ma non è cosi \\
La divisione non è una decisione che si assume \t è una situazione storica che si definisce a seguito di alcune azioni prese da una parte e dall'altra \\
La spartizione dell europa in zone di influenza non viene decisa \t ma la spartizione viene decisa dai fatti \\
La conferenza è anche l'ultima a cui partecipa Roosevelt \t era gia moribondo, e si comporta con una certa accondiscendenza con Stalin (forse anche per questo) \\
Il suo successore è Truman, e partecipera alla conferenza di Postdam

\ss[Conferenza di Postdam]
Anche questa nel 45 ma dopo Ialta \\
Qua solo il giappone è in guerra con USA \\
Truman ha atteggiamento intransigente nei confronti di Stalin \t questo cambia assetto \\
Sara cosi rigido da creare uno stallo nelle trattative \t e grazie a Churcill si sblocca questo momento difficile \\
Si decide:
\bi
\item di non smembrare la germania ma solo zone influenza
\item di riconoscere alla polonia i territori che aveva preso la germania
\item stalin si tiene una parte di polonia orientale e una parte di prussia orientale
\ei
Importante è pero che si è rotta solidarieta tra nemici del fascismo \t si è concluso l'impegno comune contro il fascismo e nazismo \\
Le differenze ideologiche e politiche emergono \t l'alleanza precedente si sta incrinando \\
Qua si pongono quindi le premesse della guerra fredda \t Truman nel 47 terrà un celebre discorso, in cui definise la Dottrina Truman, che è fortemente antibolscevico e anticomunista \\
USA si impegna a combattere ogni forma di comunismo nel mondo e a sostenere tutte forze anticomuniste in altri paesi \\
Non vuole definire una dottrina, definisce linea della sua presidenza \t ma alla storia passa come una dottrina \\
Gli USA sostiene anche economicamente europa che si pone nella stessa posizione americana

\v

Truman fa organizzare quindi il piano Marshall (che è il segretario dello stato) \t Truman lo propone a tutto il mondo, anche all'est \t fa la parte del buono, e URSS rifiuta e anche l'est europeo \\
Piano M. è anche un tentativo un po imperialista \\
In EU. occidentale arrivano 14 miliardi di dollari \t e sraa dal 48 al 51 \\
URSS risponde con un discorso di Stalin, in cui condanna questo imperialismo americano \t e propone il COMECON, che è la risposta sovietica al piano Marshall \\
Al COMECON partecipano romania, ungheria, polonia, cecoslovacchia, albania, bulgaria \t l'obbiettivo è quello di creare alleanza economica tra paesi comunisti

\v

Il 49 è un anno particolare \t nel 49 nasce infatti anche la repubblica popolare cinese \t è l'affermazione di Mao Tsetung in Cina \t che era il leader dei comunisti, che ha la meglio sui nazionalisti \t e Mao fonda la repubblica popolare (che è comunista) \\
Cina si allea (almeno all'inizio) con URSS \t ma dura poco: il comunismo cinese è molto diverso da quello sovietico \t è incentrato sul mondo contadino (no Marx) \\
I contandini si organizzano in comuni popolari \t non è il comunismo proletario e industriale stalinista \\
Stalin ha anche l'idea di stalinizzare i paesi comunisti del mondo \t vuole asservire, e per questo alleanza si rompe (e anche crisi con Tito in Jugoslavia)

\v

Il piano Marshall funziona, l'economia europea si riprende \t si rafforzano legami con USA e EU occidentale \\
Nell'aprile del 49 viene firmato il patto atlantico \t che è un accordo tra tantissimi paesi \t ed è la premessa della NATO, che è un alleanza militare che nasce a difesa del mondo libero \\
NATO viene messa anche sotto la protezione della bomba atomica americana \\
L'URSS anche lei risponde \t nel 55 \t con il patto di Varsavia = alleanza militare dei paesi comunisti, che sono tutti quei paesi in cui stalin ha imposto il comunismo (anche se veniva rifiutato dalla popolazione, come in Ungheria, Cecoslovacchia)

\v

Vengono compiute quindi azioni per cui si creano due blocchi \t ma non c'è una decisione a priori \\
La spartizione del mondo in due zone di influenza si impone come la soluzione + semplice per evitare scontri \t ognuno ha area di influenza \\
Quando eretto muro di berlino \t non c'è grande opposizione \t perchè è il segno di come questa spartizione funzionasse, e il muro di berlino la segna fisicamente \t impediva scontro diretto (anche se poi ci sono stati conflitti, pero locali, che hanno rappresentato continuo scontro ma mai diretto) \\
Nasce cosi la guerra fredda \t è una guerra mai combattuta e non armata tra le due potenze in modo diretto \t si basa sull'equilibrio generato dal terrore \\
L'equilibrio non è garantito da un accordo, ma dal terrore di una nuova guerra \t che poteva essere catastrofica = grande tensione, anche perche URSS si dota della bomba atomica (nel 48 la fa esplodere, e nel 53 ottiene la bomba a idrogeno)

\ss[Movimento dei paesi non allineati]
Non tutti i paesi sono in questo mondo spaccato \t USA e URSS si mostrano come emblema di visoni diverse \t ma alcuni dicono di no, e crea movimento dei "Paesi non allineati" \\
I paesi che fanno parte di questo gruppo sono moltissimi, ma sono per esempio tutti i paesi dell'africa, l'india etc. \t dei paesi che non hanno peso ne politico ne militare, che non possono competere \\
Questi paesi scelgono di stare fuori dalla logica di questi blocchi contrapposti \t ma non per ideologia, ma perche si vuole essere autonomi da queste logiche \t es. Jugoslavia ne fa parte ma è comunista \\
Questo movimento è anche neutralista \\
Movimento prende vita nel 55 alla conferenza di Bandung \t ci sono 29 paesi asiatici e africani, e viene promossa dall'India e dalla Jugoslavia \\
Per la prima volta si parla di neutralita e di non allineamento \\
Il movimento vero e proprio nasce pero nel 61 a Belgrado \t la Jugoslavia con Egitto e India si fa promotrice \t paesi decidono di non schierarsi ne con URSS ne con USA \\
A Lavana, ultima conferenza, ci saranno 89 paesi ma poco peso

\ss[Tito]
Jugoslavia si libera dalle forze fasciste in totale autonomia \t quindi vuole autonomia, e Tito fonda paese comunista \t Stalin vuole entrare \\
Tito si scontra con Bosca nel 47 \t Stalin cercherà di far insorgere il popolo contro Tito, ma popolo non ci crede \t e alla fine c'è rottura tra i due paesi \\
Jugoslavia quindi si allontana ancora di piu da modello sovietico, e si avvicina a paesi che volevano rimanere autonomi

\v

Da qua in avanti ci saranno molte tensioni e scontri \t che non porteranno a guerra mondiale, ma si creerà grande paura

\ss[Berlino]
Prima crisi della guerra fredda \t Berlino viene divisa in 4 zone di influenza, che sono le stesse per quelle della germania \\
Berlino pero è nella parte sovietica della germania \\
Succede quindi che nel giugno del 48 i sovietici chiudono la citta, vogliono impedirne il rifornimento \t volevano mettere in crisi la citta per prenderla \\
Gli USA rispondono con un ponte aereo, che rifornisce parte occidentale di Berlino \t forzano blocco per un anno, e URSS non risponde \t è tutta una prova di forza \\
Momento di grande tensione, ma nel maggio del 49 URSS toglie blocco \t prima grande crisi anche prima della nascita della NATO, e a ridosso della guerra
\end{document}
